\section{Historians: how doubt became a calculation}
\label{sec:historians}

Figure~\ref{fig:timeline} lays out the lineage this section reconstructs. The
five core readings sit inside four overlapping eras: the rejection rules of
observatory astronomy, the mathematical statistics of outliers, the rise of
operational network quality assurance, and the streaming era in which
statistics and software finally met the sensors.

\begin{figure}[t]
  \centering
  \includegraphics[width=0.98\textwidth]{fig_timeline.png}
  \caption{The lineage traced in this report. Marker shading groups the
  fourteen milestones into four eras; the five core readings are marked in
  the labels.}
  \label{fig:timeline}
\end{figure}

\subsection{Observatories and the first rejection rules}

The problem of the doubtful observation is as old as precision measurement,
but it became a mathematical problem only after least squares gave residuals
a definite meaning. By the 1840s an observatory reduction was an industrial
process: human computers turned raw transit timings into positions through
hours of hand arithmetic, and every astronomer knew the discomfort of a
residual too large to like but too small to blame confidently on a slip of
the pen. Practice was frankly personal. Observers struck out values on
judgement, or kept everything and let the average absorb the damage, and two
respectable astronomers could defend opposite habits with equal conviction.

Benjamin Peirce, professor of mathematics at Harvard and one of the central
figures of early American science, found that state of affairs intolerable
and in 1852 published the first general criterion for rejecting doubtful
observations \citep{peirce1852}. His principle was characteristically
ambitious: reject a subset of observations when the probability of the whole
system of errors, computed with the suspects excluded and the rejection
itself accounted for, exceeds the probability computed with them retained.
The criterion handled several outliers at once and required an iterative
calculation that pushed the limits of hand computation. It was quickly taken
up around the United States Coast Survey, the great measurement bureaucracy
of the day, and Benjamin Gould published a worked guide to make it usable by
practical computers \citep{gould1855}. It also drew immediate fire. George
Airy, the Astronomer Royal, argued that no formal rule could substitute for
the observer's knowledge of the circumstances of each measurement, and that
mechanical rejection was a licence to tidy away inconvenient evidence. The
whole modern argument about automated flagging, judgement versus procedure,
transparency versus convenience, was thus present at the birth.

William Chauvenet, who had built the mathematics curriculum of the new United
States Naval Academy and would later lead Washington University in St.\
Louis, took the pragmatic path. In an appendix to his astronomy manual of
1863 he proposed the rule that still carries his name: reject an observation
if, under the fitted normal law, the expected number of deviations as large
as the one observed is less than one half in a sample of the given size
\citep{chauvenet1863}. Chauvenet's criterion is crude, its implicit
significance level drifts with sample size, and he offered it modestly as a
convenience. It won anyway, for a century, because a computer working by
candlelight could apply it in one line. The lesson, that adoption follows
usability at least as much as principle, will recur in every era of this
story.

\subsection{From rejection rules to mathematical statistics}

The founders of modern statistics inherited the rejection question and
changed its form. Once distributions of sample statistics could be derived
exactly, the natural move was to treat the largest residual as a test
statistic and ask for its null distribution. William Thompson proposed
studentising the extreme deviate in 1935 \citep{thompson1935}, and a year
later Egon Pearson and C.\ Chandra Sekar published the analysis that gave
the subject its first real depth \citep{pearson1936}. They showed that a test built on the most
extreme observation can fail completely when two or more outliers stand
together, each hiding the other by inflating the sample scale. They named the
phenomenon masking. It is a beautiful early example of a failure mode
discovered on paper decades before practitioners would routinely trip over it
in data, and Section~\ref{sec:experimentalists} reproduces it numerically.

The person who turned this thread into a working toolkit spent his career a
long way from any observatory. Frank Ephraim Grubbs, born in Alabama in 1913,
joined the Army's Ballistic Research Laboratory at Aberdeen Proving Ground
before the Second World War and stayed for his whole working life, completing
a doctorate at Michigan in 1949 under Cecil Craig with a dissertation on
sample criteria for outlying observations \citep{grubbs1950}. The setting
explains the statistics. Ballistic testing destroys its own sample: every
fired round is money and time, samples are small by necessity, and one
aberrant round in a lot of twenty can decide an acceptance test. Small
samples, expensive observations and consequential decisions are exactly the
regime where a principled outlier test earns its keep. Grubbs became one of
the Army's most decorated statisticians, receiving the American Statistical
Association's Wilks Memorial Award in 1964.

His 1969 paper in \work{Technometrics}, the first core reading, is a work of
synthesis rather than novelty, and that is precisely its historical
importance \citep{grubbs1969}. \work{Technometrics} had been founded ten
years earlier, jointly by the statistics and quality control communities, to
serve statistics in the physical sciences and industry, and Grubbs wrote for
its readership: engineers who needed recipes, tables and worked examples
rather than lemmas. The paper collects the extreme studentised deviate tests,
his own two-outlier statistics, sample kurtosis, and the accumulated critical
value tables into a single guide, organised by the practical question being
asked. It won the journal's Youden and Wilcoxon prizes in 1969, and it
remains, as Section~\ref{sec:futurists} shows, one of the steadily most cited
statistics papers of its generation. A year earlier Frank Anscombe had
supplied the era's other lasting idea: rejection rules should be judged like
insurance policies, by the premium you pay in good observations wrongly
discarded against the protection you buy when something is genuinely wrong
\citep{anscombe1960}. Every false alarm budget in every modern monitoring
system is Anscombe's premium under a new name.

\subsection{The robust turn}

Grubbs' toolkit rests on the sample mean and standard deviation, and by the
1960s the statistical avant garde had turned against both. Tukey's work on
contaminated distributions showed how little contamination it takes to wreck
the efficiency of classical estimators, Huber built a rigorous theory of
estimators that tolerate it, and Hampel added the vocabulary of influence
functions and breakdown points that made robustness precise
\citep{hampel1974}. Out of that programme came a modest pair of estimators
with maximal resistance: the sample median for location, and the median of
absolute deviations from the median, MAD, for scale, both able to withstand
corruption of just under half the sample. The philosophical shift mattered as
much as the mathematics. Where the classical school asked which observations
to reject before estimating, the robust school preferred to accommodate:
estimate in a way that suspect observations cannot dominate, and rejection
becomes a by-product rather than a prerequisite.

The second core reading is the bridge that carried this material from the
research literature to the workbench. The American Society for Quality
Control commissioned a series of short booklets putting usable statistics in
the hands of quality engineers, and volume sixteen, \work{How to Detect and
Handle Outliers}, was written by Boris Iglewicz of Temple University, a
quality methods specialist, and David Hoaglin, one of the editors of the
robust analysis handbooks that codified the Tukey school
\citep{iglewicz1993}. The booklet walks an engineer from boxplots through
formal tests, but its lasting export is three lines long: replace the mean by
the median, the standard deviation by a rescaled MAD, and flag any
observation whose modified z-score exceeds 3.5 in absolute value. The
recommendation was calibrated by simulation for batches of independent
measurements. Neither author was thinking about a sensor reporting every five
minutes, a fact that will matter a great deal in
Section~\ref{sec:experimentalists}.

\subsection{A flood, a network, and quality assurance as a discipline}

The operational era begins, by common consent of its own founders, with a
flood. Over the Memorial Day weekend of 1984 a stalled storm dropped
torrential rain on Tulsa in the middle of the night, killing fourteen people
and causing damage the sparse observing network had given little warning of.
Oklahoma's meteorologists concluded that the state simply could not see its
own weather, and an unusual alliance of the state's two rival universities,
meteorology at the University of Oklahoma and agriculture at Oklahoma State,
spent the next decade building the answer: the Oklahoma Mesonet, a network of
over a hundred automated stations, roughly one per county, commissioned at
the start of 1994 \citep{brock1995}. Solar powered stations logged more than
a dozen quantities every five to thirty minutes and relayed them through the
state law enforcement telecommunications system, giving Oklahoma a live,
state-wide picture of the mesoscale atmosphere that no one had possessed
before.

It also gave Oklahoma a problem no one had faced at that scale: hundreds of
thousands of unattended readings a day, promised to users within minutes, on
instruments sitting in fields. The Mesonet's answer, documented in the third
core reading \citep{shafer2000}, was to treat data quality as an engineering
discipline with its own staff, budget and paper trail. Instruments were
calibrated in a dedicated laboratory before deployment, rotated on schedule,
and compared in the field against travelling reference sensors. Incoming data
passed a battery of automated tests, range, step, persistence,
like-instrument and spatial comparisons, whose verdicts were compressed into
a graded flag rather than a deletion. And, crucially, the automated verdicts
were treated as advice to a human quality assurance meteorologist, who
reviewed daily summaries, opened trouble tickets and fed confirmed faults
back to the technicians. Nothing in the statistical content was new; the
architecture was. A parallel and mathematically deeper tradition already
existed in numerical weather prediction, where Gandin's complex quality
control cross-examined every observation against an optimal interpolation of
its neighbours \citep{gandin1988}, but that machinery assumed the dense
infrastructure of a forecasting centre. The Mesonet showed what a lean
institution could sustain indefinitely, and its blueprint, simple tests plus
graded flags plus human review, became the de facto standard for surface
networks worldwide \citep{fiebrich2010, estevez2011}.

\subsection{The sensor deluge and the return of the statisticians}

Through the 2000s the cost of sensing collapsed. Ecology, hydrology and
limnology bolted sondes, sap flow sensors and micrometeorological stations
onto research sites that had been sampled monthly by hand, grassroots
networks such as the lake observatory network GLEON federated the streams,
and in the United States the National Ecological Observatory Network was
being designed to run thousands of instruments for thirty years
\citep{taylor2013}. Ecological data volumes grew a thousandfold while
ecological data skills did not, and by the early 2010s the community's
methods papers carried an unmistakable tone of alarm. The fourth core
reading is the framework that answered it. Written by John Campbell, Lindsey
Rustad and colleagues rooted in the long-term forest research community, it
translated the Mesonet's institutional wisdom into the vocabulary of
streaming ecological informatics: quality assurance as everything designed
before the data flow, quality control as the checks applied after, a
taxonomy of sensor faults distinct from scientifically interesting
anomalies, graded flags, versioned archives that never silently overwrite,
and middleware that carries all of it \citep{campbell2013}. Its title,
\work{Quantity is Nothing without Quality}, became the field's slogan.

The fifth reading closes the circle by bringing professional statisticians
back to a problem astronomers had handed them 160 years earlier. In
Queensland, the agency monitoring river loads into the Great Barrier Reef
lagoon was drowning in high-frequency turbidity and conductivity data whose
manual validation could not keep pace, and it partnered with the Australian
centre of excellence for mathematical statistics to automate the triage. The
resulting framework paper \citep{leigh2019}, with forecasting, spatial
statistics and Bayesian specialists among its authors, defined anomaly types
from the sensors' actual failure modes, compared regression-based,
feature-based and rule-based detectors on expert-labelled series, and
reported performance honestly by type. Its companion papers put the methods
into open source software \citep{talagala2019}. After 170 years the wheel
had come full circle: the discipline that began with astronomers arguing
over one doubtful transit now fields statisticians, domain scientists and
software engineers jointly deciding, at five minute resolution, what a river
is allowed to do.

One question the Historian role poses deserves a systematic answer: what
could the computers of each era actually do? Table~\ref{tab:computing}
summarises the machinery available to each of the core readings, because in
this subject the machinery is never incidental. Peirce's criterion lost to
Chauvenet's partly because iteration was expensive when arithmetic was
human; Grubbs' tables were the interface between exact theory and engineers
who had no computer at all; the Mesonet's five minute cadence was set as
much by radio bandwidth and datalogger memory as by meteorology; and the
current era's appetite for learned methods rides on hardware that makes
refitting a model per station per day unremarkable. Methods succeed when
they fit their machines.

\begin{table}[t]
  \centering
  \small
  \caption{The computing context of each core reading.}
  \label{tab:computing}
  \begin{tabular}{@{}lp{4.0cm}p{6.9cm}@{}}
    \toprule
    Work & Computing of its time & Consequence for the method \\
    \midrule
    \citet{grubbs1969} & institutional mainframes; practitioners with slide
      rules and printed tables & exact null distributions tabulated once,
      centrally, and used everywhere; tests designed around small $n$ \\
    \citet{iglewicz1993} & desktop PCs, statistical packages on every
      engineer's desk & robust estimators affordable interactively; a
      three-line rule could realistically be applied to every batch \\
    \citet{shafer2000} & solar-powered dataloggers, radio telemetry,
      central Unix servers & automated batteries run centrally on schedule;
      test complexity capped by the daily batch and by telemetry bandwidth \\
    \citet{campbell2013} & wireless sensor networks, middleware, relational
      archives & streaming checks at ingest; versioned archives and flag
      columns become cheap, so nothing need be deleted \\
    \citet{leigh2019} & open source R and Python, cloud computing &
      per-site model fitting and simulation-based thresholds affordable;
      code shipped with the paper \\
    \bottomrule
  \end{tabular}
\end{table}

Reading the five works as history, one competing methodology deserves its own
epitaph. At every stage there was a purer alternative that lost on
operational grounds: Peirce's principled criterion lost to Chauvenet's one
liner, Gandin's complex quality control stayed inside forecasting centres
while mesonets ran simple rule batteries, and the robust school's counsel to
accommodate outliers rather than hunt them was largely ignored by monitoring
practice, which needs named, actionable faults rather than quietly
downweighted ones. The survivors are not the most elegant methods but the
ones an institution could afford to run, explain and staff. That selection
pressure, more than any theorem, shaped the systems we operate today.
